% (User input window)

\subsubsection{有界切片上的连通性、可逆化与残差几何}\label{subsubsec:xi-bounded-slice-connectivity-residual-geometry}
在定义 \ref{def:pom-bounded-slice} 的任意固定有界复杂度切片 \(\mathbf{PW}^{\ast}_{\le}\) 上，保守语义函子 \(\mathsf{Sem}\)（定义 \ref{def:pom-sem-functor}；定理 \ref{thm:pom-sem-conservative}）将实现差异收束为自然同构类；
同时，交换失败残差在可组合口径下被异常上同调类 \([\mathrm{Anom}]\) 承载（\S\ref{subsec:pom-anom-cohomology}）。
在该框架内，“连通性增长”可被严格拆解为：可逆化导致的同构类几何、残差承载的循环自由度、以及严格化所受的一致性障碍。

\begin{conclusion}[审计可逆的正交因子分解与连通性正规形的唯一性]\label{con:xi-connectivity-orthogonal-factorization-normal-form}
设 \(\rho:\mathrm{Ob}(\mathbf{PW}^{\ast}_{\le})\to\NN\) 为分辨率等级函数，并在 \(\mathbf{PW}^{\ast}_{\le}\) 上固定一套可审计的残差核算口径，使得：
\begin{enumerate}
  \item 对任意可复合 \(2\)-态射证书，其残差累积满足可复核的可加性（定义 \ref{def:pom-anom-accumulation}；命题 \ref{prop:pom-anom-additive-up-to-boundary}）；
  \item \(\mathsf{Sem}\) 在该切片上保守，从而反射自然同构（定理 \ref{thm:pom-sem-conservative}）。
\end{enumerate}
定义两类 \(1\)-态射：\(\mathcal L\) 为分辨率严格不降（\(\rho\) 单调不减）的提升类，\(\mathcal P\) 为分辨率严格不升（\(\rho\) 单调不增）的读出类。
则对任意 \(1\)-态射 \(w\) 存在分解
\[
w\simeq p\circ \ell,\qquad \ell\in\mathcal L,\ p\in\mathcal P,
\]
并且该分解在 \(2\)-同构意义下唯一。
更强地，\((\mathcal L,\mathcal P)\) 构成一个稳定正交分解系统：任意“读出--提升”交换方块存在填充且填充至唯一 \(2\)-同构。
因此，任意由生成门给出的概念连通路径在该切片内都存在唯一（至 \(2\)-同构）的规范形代表，并且在残差与复杂度意义下同时达到最小。
\end{conclusion}

\begin{conclusion}[语义等价类的图论刻画：强连通分量与可逆闭包]\label{con:xi-connectivity-scc-invertible-closure}
在固定有界切片 \(\mathbf{PW}^{\ast}_{\le}\) 上，将所有被 \(\mathsf{Sem}\) 识别为可逆的 \(1\)-态射形式地可逆化，得到局部化群胚 \(\mathcal{G}_{\le}\)。
取生成 \(1\)-态射得到有向多重图 \(G_{\le}\)（顶点为对象，边为生成态射）。
若在该切片上可审计重写足以判定 \(2\)-同构（从而“同一 \(\mathsf{Sem}\) 输出 \(\Rightarrow\) 自然同构”可被有限证书核对；定理 \ref{thm:pom-sem-conservative}），则 \(G_{\le}\) 的强连通分量分解与 \(\mathcal{G}_{\le}\) 的同构类分解一致。
等价地，两对象在 \(G_{\le}\) 中互达，当且仅当在 \(\mathbf{PW}^{\ast}_{\le}\) 中存在一对互为逆的可审计投影词（至 \(2\)-同构）。
因此，在不改变语义切分的前提下，加边只能改变同一语义类内的导航直径与分量间凝聚 DAG 的可达代价，而不能产生新的语义同构类。
\end{conclusion}

\begin{conclusion}[分辨率过滤下的持久连通谱：稳定概念的条形码刻画]\label{con:xi-connectivity-persistent-barcode}
令预算/分辨率上界 \(B\mapsto \mathbf{PW}^{\ast}_{\le B}\) 给出一族单调嵌入的有界切片，并令 \(\mathcal{G}_{\le B}\) 为其在 \(\mathsf{Sem}\)-可逆态射上的局部化群胚。
记 \(\pi_0(\mathcal{G}_{\le B})\) 为同构类集合，则包含函子诱导出直系系统
\(
\pi_0(\mathcal{G}_{\le B})\to \pi_0(\mathcal{G}_{\le B'})
\)（\(B\le B'\)）。
若对每个 \(B\) 该集合有限，且类的合并只能由可审计证书触发（不会发生不可追踪坍塌），则该过滤唯一确定一组区间多重集 \(\mathcal I\)（持久连通谱/条形码），满足：
\begin{enumerate}
  \item 无限区间与稳定语义类一一对应：存在 \(B_0\) 使得对所有 \(B\ge B_0\) 该类在 \(\pi_0(\mathcal{G}_{\le B})\) 中保持不变；\label{item:xi-persist-infinite}
  \item 有限区间对应瞬时同构或瞬时概念，其出现与消失可由残差与证书定位；\label{item:xi-persist-finite}
  \item 稳定概念可在逆极限意义下刻画为 \(\varprojlim_B \mathcal{G}_{\le B}\) 中的非 \(\mathsf{NULL}\) 可寻址对象族（参见定理 \ref{thm:xi-time-consistent-inverse-limit-readout} 的同型逆系统论证）。\label{item:xi-persist-invlim}
\end{enumerate}
于是，连通性随预算增长的可观察变化可被重述为：在保持 \ref{item:xi-persist-infinite} 所述无限条不变的前提下，允许有限条形码发生可审计迁移，从而把体系增长约束为对持久谱的受控扰动。
\end{conclusion}

\begin{conclusion}[连通性增长的不可逃代价：循环秩由异常维数下界控制]\label{con:xi-connectivity-betti-lower-bound-by-anom-rank}
令 \(G_{\le B}\) 为预算 \(B\) 下的生成图（可取为结论 \ref{con:xi-connectivity-scc-invertible-closure} 的 \(G_{\le}\) 在阈值 \(B\) 下的子图），并令 \(\beta_1(G_{\le B})\) 表示其无向化后的第一 Betti 数。
取阿贝尔群 \(\mathrm{Anom}_{\le B}\) 表示在切片 \(\mathbf{PW}^{\ast}_{\le B}\) 上可核对的异常自由度（例如由定义 \ref{def:pom-anom-accumulation} 的残差坐标生成），并设存在闭路残差映射
\(
\alpha_B:Z_1(G_{\le B})\to \mathrm{Anom}_{\le B}
\)
将每个闭路送到其规范化审计残差。
若零残差闭路必可由 \(2\)-同构填充消去（即残差外置最小性），则 \(\alpha_B\) 在同调意义上诱导单射
\[
H_1(G_{\le B};\ZZ)\hookrightarrow \mathrm{Anom}_{\le B},
\]
从而得到结构性下界
\[
\beta_1(G_{\le B})\le \mathrm{rank}_{\ZZ}(\mathrm{Anom}_{\le B}).
\]
因此，若试图通过加边提高连通性密度并同时保持语义不坍塌，则异常自由度必须同比例增长以承载新增的独立循环；否则新增循环将被迫塌缩为可逆冗余，不产生新的可用结构。
\end{conclusion}

\begin{conclusion}[审计保持严格化的三阶障碍：连通性闭合的高阶一致性判据]\label{con:xi-connectivity-third-order-strictification-obstruction}
在 \(\mathbf{PW}^{\ast}_{\le B}\) 的双范畴结构中，连通性不仅依赖 \(1\)-态射可达性，还依赖 \(2\)-态射的相容性（不同括号化/不同重写路径的审计一致）。
由三重复合的括号化差可定义一个取值于 \(\mathrm{Anom}_{\le B}\) 的三阶一致性障碍类
\[
\omega_B\in H^3(G_{\le B};\mathrm{Anom}_{\le B}),
\]
可视为审计保持的结合子缺陷的同调类。
则：存在一个在 \(2\)-范畴意义下与 \(\mathbf{PW}^{\ast}_{\le B}\) 等价、且保持 \(\mathsf{Sem}\) 与残差核算的严格 \(2\)-范畴模型，当且仅当 \(\omega_B=0\)。
因此，将任意多段连通路径规约到唯一规范形且审计一致的“可机器导航闭合”，受制于一个可计算的三阶一致性障碍；当该障碍为零时，连通性可在不增添额外寄存器自由度的前提下严格闭合。
\end{conclusion}

\begin{conclusion}[最小知识图谱骨架的存在与计数：连通可用性的线性规模实现]\label{con:xi-connectivity-minimal-skeleton-linear-size}
设预算 \(B\) 下对象数为 \(n_B:=|\mathrm{Ob}(\mathbf{PW}^{\ast}_{\le B})|\)，并假设 \(\mathrm{Anom}_{\le B}\) 为自由阿贝尔群，且闭路残差映射 \(\alpha_B\) 在同调上满射（即异常维度在切片上可由闭路证书完全激发）。
则存在一个子图 \(H_B\subseteq G_{\le B}\)（审计等价骨架）满足：
\begin{enumerate}
  \item \(H_B\) 与 \(G_{\le B}\) 诱导相同的局部化群胚（同构类与可达关系一致）；\label{item:xi-skel-groupoid}
  \item \(\alpha_B\) 在 \(H_B\) 上仍保持同调满射（异常自由度不丢失）；\label{item:xi-skel-anom}
  \item 边数达到不可改进下界
  \[
  |E(H_B)| = n_B - c_B + \mathrm{rank}_{\ZZ}(\mathrm{Anom}_{\le B}),
  \]
  其中 \(c_B\) 为凝聚图的弱连通分量数。\label{item:xi-skel-size}
\end{enumerate}
因此，在保守语义与异常审计充分激发的前提下，任何可用的概念连通性都可压缩到线性规模骨架，且骨架大小由“对象数 \(+\) 必要异常维度”精确决定；若异常维度不随连通性增长而增长，则新增边最终仅形成可删去的冗余而不提升可用结构覆盖率。
\end{conclusion}

